\subsection{Sim-to-Real Robustness}
\label{sec:robustness}
\label{sec:clustered_curriculum}

A test-time sweep checks whether the architectural advantage holds outside the
training distribution. Neither policy is retrained. Four perturbation axes are
varied: GPS position noise ($\sigma\in\{0,5,10,20\}$ grid units), path-loss
exponent mismatch ($n\in\{3.4,3.8,4.2,4.6\}$), log-normal shadowing variance
($\sigma_{\mathrm{sh}}\in\{4,6,8,10\}$\,dB), and wind drift
($w\in\{0,0.5,1.0,2.0\}$ cells/step). Each (agent, axis, magnitude) cell uses
10 seeds on $(200{\times}200, N{=}20)$, deterministic $\arg\max$. Per-seed data
is in \texttt{baseline\_results/sim\_to\_real/robustness\_raw.csv}; NDR curves
are in Figure~\ref{fig:robustness_summary} (Appendix~\ref{app:extended_results}).

On the three working axes neither policy shows a perturbation cliff. The
Relational-RL policy holds at $100\%$ NDR at every magnitude; the DQN drops
by at most $\approx 1$ percentage point. The path-loss panel is a
\emph{perturbation no-op}: the Two-Ray model in
\texttt{iot\_sensors.py:\allowbreak calculate\_rssi()} uses hardcoded slope
constants, making \texttt{path\_loss\_exponent} inert. Replacing it with a
parametric log-distance model is high-priority future work
(Section~\ref{sec:future_work}). The architectural advantage holds on three of
the four axes. Per-axis detail and the shadowing non-monotone reward observation
are in Appendix~\ref{app:robustness_extended}.
